% Provenance (never printed): see sources/outline_sources.md (Year 2: eccentricity, obliquity and
% precession; insolation computed from spherical astronomy; the Quaternary cycles; carbon dioxide
% and feedbacks). Orbital elements and tabulated insolation of Berger and Loutre (1991) via NOAA
% NCEI (data/derived/bachelor-2-orbit91, fetch_data Book 4 block); periods after Laskar et al.
% (2004) and the NASA Milankovitch page; history after the NASA Earth Observatory feature (2000);
% terminations after Barker and Knorr (2021); the mid-Pleistocene transition after Clark et al.
% (2006) and Koehler and van de Wal (2020). Figures from figdata/bachelor-2/orbital.py; values of
% the exercises and problem computed in the Book 4 scratch calculator calc/ch25.py.
\chapter{Gaya Paksa Orbital dan Zaman Es}\label{ch:b2:orbital-forcing}

Selama sejuta tahun terakhir bentuk orbit Bumi telah mengubah sinar Matahari yang diterima
seluruh planet dalam setahun kurang dari dua per seribu: antara \num{340.00} dan
\qty{340.54}{W.m^{-2}} di puncak atmosfer. Selama sejuta tahun yang sama, penyinaran harian rata-rata
pertengahan Juli pada \ang{65} LU telah berayun antara sekitar 388 dan \qty{500}{W.m^{-2}}, lebih dari
seperempatnya. Zaman es mengikuti angka kedua, bukan angka pertama. Bab ini menghitung keduanya
dari unsur-unsur orbit Bumi, dan bertanya bagaimana sinar Matahari pada satu lintang dalam satu musim dapat
mengatur laju pertumbuhan dan peluruhan lembaran es yang menutupi benua.
% ledger: data:b4-orbital, berger91:readme

\begin{recall}
\textbf{Dari buku-buku fisika}: elips, sumbu semi-mayor $a$ dan eksentrisitasnya $e$, hukum-hukum
Kepler. \textbf{Dari volume Tahun 1}: ekliptika dan kemiringan ekliptika $\varepsilon$, deklinasi
Matahari $\delta$ dengan $\sin\delta = \sin\varepsilon\sin\lambda$ ($\lambda$ bujur sejati
Matahari, yang dihitung di sepanjang ekliptika dari titik ekuinoks musim semi), tahun tropis, dan insolasi
harian rata-rata
\[
  \bar Q = \frac{S_0}{\pi}\left(\frac{a}{r}\right)^{2}\bigl(h_0\sin\varphi\sin\delta +
  \cos\varphi\cos\delta\sin h_0\bigr),\qquad \cos h_0 = -\tan\varphi\tan\delta ,
\]
serta rata-rata global dan tahunannya, $S_0/(4\sqrt{1 - e^2})$. \textbf{Dari volume-volume sekolah}:
siklus glasial dan Maksimum Glasial Terakhir; gaya radiatif. \textbf{Dari bab-bab sebelumnya}:
\omterm{def:b2:precambrian-earth:ice-albedo}{umpan balik es--albedo} (\cref{def:b2:precambrian-earth:ice-albedo}); tumpukan bentik, rekaman EPICA
Dome C, dan \omterm{def:b2:climate-proxies:age-model}{model umur} (\cref{ch:b2:climate-proxies}).
\end{recall}
% ledger: earth:TSI, earth:obliquity, earth:ecc

\section{Orbit dan ketiga unsurnya}\label{sec:b2:orbital-forcing:elements}

Orbit Bumi adalah elips dengan Matahari di salah satu fokusnya. Planet-planet, terutama Jupiter dan Saturnus,
perlahan-lahan mengubah bentuknya dan memutarnya di dalam bidangnya; tarikan Matahari dan Bulan pada tonjolan
khatulistiwa membuat sumbu rotasi bergoyang seperti gasing. Tiga besaran menggambarkan apa yang penting bagi
sinar Matahari: seberapa lonjong elips itu, seberapa miring sumbunya, dan di mana di sepanjang orbit
musim-musim jatuh.
% ledger: nasa:milank-causes

\begin{definition}[Bujur perihelion; presesi iklim]\label{def:b2:orbital-forcing:elements}
Yang disebut \emph{bujur perihelion}\index{bujur perihelion} $\varpi$ adalah sudut, yang diukur di
bidang orbit dari arah titik ekuinoks musim semi, antara arah itu dan
arah perihelion dilihat dari Matahari; sudut ini dihitung dari ekuinoks yang bergerak, sehingga
berubah karena goyangan sumbu maupun perputaran orbit. Yang disebut \emph{presesi
iklim}\index{presesi iklim} adalah besaran $e\sin\varpi$, yang juga disebut indeks
presesi.
\end{definition}
% ledger: berger91:readme

\begin{omfigure}
\centering
\begin{tikzpicture}[scale=0.82, font=\scriptsize, >=Stealth]
\draw[thick, rotate=101.37] (-0.72,0) ellipse (2.4 and 2.2895);
\draw[densely dashed, black!60] (-0.331,1.647) -- (0.615,-3.059);
\draw[->, black!70] (0,0) -- (3.9,0);
\node[anchor=south, align=center] at (3.75,0.06) {$\gamma$, ke arah\\titik ekuinoks\\musim semi};
\draw[->, omDef] (0.75,0) arc[start angle=0, end angle=101.37, radius=0.75];
\node[omDef] at (0.78,0.62) {$\varpi$};
\fill[ocRed] (0,0) circle (0.13);
\node[anchor=north west] at (0.05,-0.08) {Matahari};
\foreach \x/\y in {-2.062/0, 0/-3.094, 2.321/0, 0/1.688} {\fill[ocBlue] (\x,\y) circle (0.08);}
\node[anchor=east, align=right] at (-2.15,0) {ekuinoks\\Maret};
\node[anchor=north west, align=left] at (2.38,-0.08) {ekuinoks\\September};
\node[anchor=west] at (0.1,1.92) {solstis Desember};
\node[anchor=east] at (-0.12,-3.25) {solstis Juni};
\fill[black] (-0.331,1.647) circle (0.06);
\node[anchor=east] at (-0.42,1.72) {perihelion};
\fill[black] (0.615,-3.059) circle (0.06);
\node[anchor=west] at (0.7,-3.15) {aphelion};
\draw[->, thick, black!60] (-1.9,-1.75) arc[start angle=220, end angle=245, radius=2.6];
\end{tikzpicture}
\par\smallskip{\small Orbit Bumi dilihat dari sisi utara, dengan eksentrisitas yang dilebih-lebihkan
menjadi 0.3 dan \omterm{def:b2:orbital-forcing:elements}{bujur perihelion} masa kini, $\varpi = \ang{101.4}$. Bumi bergerak berlawanan arah jarum jam
(panah); titik-titik menandai kedudukannya pada ekuinoks dan solstis. Perihelion jatuh beberapa
hari setelah solstis Desember, aphelion beberapa hari setelah solstis Juni. Skematis.}
\end{omfigure}
% ledger: berger91:rows, nasa:milank-detail

Dilihat dari Bumi, Matahari bergerak di sepanjang ekliptika, dan bujur sejatinya $\lambda$ adalah
sudut dari titik ekuinoks musim semi: $\lambda = \ang{90}$ pada solstis Juni, \ang{270} pada solstis
Desember. Bumi berada di perihelion ketika Matahari berada pada arah yang berlawanan dengan perihelion,
pada $\lambda = \varpi + \ang{180}$; kini $\varpi = \ang{101.37}$ dan perihelion jatuh pada
$\lambda = \ang{281.37}$, sekitar sebelas hari setelah solstis Desember, pada atau sekitar 3 Januari.
% ledger: berger91:rows, nasa:milank-detail

\begin{definition}[Parameter Milankovitch]\label{def:b2:orbital-forcing:milankovitch}
Yang disebut \emph{parameter Milankovitch}\index{parameter Milankovitch} adalah ketiga besaran
orbital yang menentukan sebaran sinar Matahari di Bumi menurut musim dan
lintang: eksentrisitas $e$, kemiringan ekliptika $\varepsilon$, dan \omterm{def:b2:orbital-forcing:elements}{presesi iklim}
$e\sin\varpi$. Yang disebut \emph{teori Milankovitch}\index{teori Milankovitch} menyatakan bahwa variasi lambat
ketiga parameter itu mengatur laju pertumbuhan dan peluruhan lembaran es Kuarter, melalui insolasi musim panas
di lintang tinggi utara.
\end{definition}
% ledger: nasa:milank-detail, eo00:milankovitch

Setiap parameter bervariasi dengan iramanya sendiri. Eksentrisitas memiliki periode utama dekat \num{95000},
\num{99000}, \num{124000}, dan \num{131000} tahun, yang bergabung menjadi siklus sekitar
\num{100000} tahun, serta periode \num{405000} tahun, yang paling stabil di antara semuanya; selama sejuta
tahun terakhir eksentrisitas berkisar dari \num{0.0046} sampai \num{0.0565}, dan kini, pada \num{0.0172}, perlahan
menurun. Kemiringan sumbu bervariasi dengan periode utama \num{40978} tahun, antara \ang{22.1} dan
\ang{24.5} selama sejuta tahun terakhir; pada \ang{23.44} kemiringan itu sedang menurun, setelah melewati maksimum terakhirnya
sekitar \num{10000} tahun yang lalu. Sumbu bergoyang mengitari satu putaran dalam sekitar \num{25770} tahun, sementara
orbit berputar dalam sekitar \num{112000} tahun; bersama-sama keduanya menggerakkan perihelion melalui musim-musim dengan
periode utama dekat \num{19000}, \num{22400}, dan \num{23700} tahun.
% ledger: la04:eccentricity, la04:obliquity, la04:precession, nasa:milank-detail, data:b4-orbital

\begin{omfigure}
\centering
% figdata/bachelor-2/orbital.py part elements
\begin{tikzpicture}
\begin{groupplot}[group style={group size=1 by 3, vertical sep=0.3cm, xlabels at=edge bottom,
  xticklabels at=edge bottom}, width=0.96\linewidth, height=3.1cm, xmin=0, xmax=1000,
  xtick={0,100,...,1000}, tick label style={font=\scriptsize}, label style={font=\small},
  scaled y ticks=false, yticklabel style={/pgf/number format/fixed,
  /pgf/number format/precision=2}]
\nextgroupplot[ylabel={$e$}, ymin=0, ymax=0.085, ytick={0,0.02,0.04,0.06}]
\addplot[omDef, thin] table[x=age_ka, y=ecc] {figdata/out/bachelor-2-orbital-elements.dat};
\node[font=\scriptsize, anchor=north west] at (rel axis cs:0.01,0.97) {a\enspace eksentrisitas};
\nextgroupplot[ylabel={$\varepsilon$ (\unit{\degree})}, ymin=21.9, ymax=25.4, ytick={22,23,24}]
\addplot[ocBlue, thin] table[x=age_ka, y=obl] {figdata/out/bachelor-2-orbital-elements.dat};
\node[font=\scriptsize, anchor=north west] at (rel axis cs:0.01,0.97) {b\enspace kemiringan sumbu};
\nextgroupplot[ylabel={$e\sin\varpi$}, ymin=-0.07, ymax=0.12, ytick={-0.05,0,0.05},
  xlabel={umur (ribuan tahun sebelum 1950)}]
\addplot[ocRed, thin] table[x=age_ka, y=prec] {figdata/out/bachelor-2-orbital-elements.dat};
\node[font=\scriptsize, anchor=north west] at (rel axis cs:0.01,0.97) {c\enspace presesi
iklim};
\end{groupplot}
\end{tikzpicture}
\par\smallskip{\small Ketiga \omterm{def:b2:orbital-forcing:milankovitch}{parameter Milankovitch} selama sejuta tahun terakhir: (a)
eksentrisitas; (b) kemiringan sumbu; (c) \omterm{def:b2:orbital-forcing:elements}{presesi iklim}, yang selubungnya adalah eksentrisitas.
Data: Berger dan Loutre (1991), NOAA NCEI Paleoclimatology, sebagaimana diambil pada 23 Oktober 2026.}
\end{omfigure}
% ledger: data:b4-orbit91

\omterm{def:b2:orbital-forcing:elements}{Presesi iklim} berayun cepat, tetapi hanya sejauh yang dibolehkan eksentrisitas: ketika orbit
hampir bulat, seperti \num{400000} tahun yang lalu, hampir tidak penting di mana perihelion jatuh.
Gambar itu menunjukkannya dengan jelas: selubung panel (c) adalah panel (a).
% ledger: data:b4-orbit91

\section{Sinar Matahari dari orbit}\label{sec:b2:orbital-forcing:insolation}

Insolasi harian rata-rata yang diingatkan di atas memerlukan satu bahan lagi: jarak Matahari pada
hari tertentu.

\begin{proposition}[Jarak dan bujur Matahari]\label{prop:b2:orbital-forcing:distance}
Bila bujur sejati Matahari adalah $\lambda$, jarak Bumi--Matahari adalah
\[
  \frac{r}{a} = \frac{1 - e^2}{1 - e\cos(\lambda - \varpi)} .
\]
Pada solstis Juni ($\lambda = \ang{90}$), sampai orde pertama dalam $e$,
$(a/r)^2 \approx 1 - 2e\sin\varpi$; pada solstis Desember, $(a/r)^2 \approx 1 + 2e\sin\varpi$.
\end{proposition}

\begin{proof}
Pada elips dengan Matahari di salah satu fokus, $r/a = (1 - e^2)/(1 + e\cos v)$, dengan anomali sejati $v$
yaitu sudut di Matahari dari perihelion ke Bumi. Bumi dilihat dari Matahari pada
arah yang berlawanan dengan arah Matahari dilihat dari Bumi, sehingga bujurnya dilihat dari Matahari adalah
$\lambda + \ang{180}$, dan $v = \lambda + \ang{180} - \varpi$, $\cos v = -\cos(\lambda - \varpi)$.
Pada $\lambda = \ang{90}$, $\cos(\lambda - \varpi) = \sin\varpi$ dan $(a/r)^2 = (1 -
e\sin\varpi)^2/(1 - e^2)^2 = 1 - 2e\sin\varpi + O(e^2)$; pada $\lambda = \ang{270}$ tandanya berubah.
\end{proof}

Karena itu \omterm{def:b2:orbital-forcing:elements}{presesi iklim} adalah indeks yang tepat: indeks ini mengukur seberapa lebih dekat ke Matahari, atau
lebih jauh, Bumi berada pada solstis Juni. Bila $e\sin\varpi$ besar dan positif, seperti masa kini dan
jauh lebih besar lagi \num{116000} tahun yang lalu, musim panas utara jatuh dekat aphelion dan lemah; bila
besar dan negatif, seperti sekitar \num{10000} tahun yang lalu, musim panas utara jatuh dekat perihelion.
% ledger: berger91:rows

\begin{example}[Matahari Juni pada \ang{65} LU, masa kini]\label{ex:b2:orbital-forcing:today}
Kini $e = 0.017236$, $\varpi = \ang{101.37}$, $\varepsilon = \ang{23.446}$, dan tabel-tabelnya memakai
$S_0 = \qty{1360}{W.m^{-2}}$. Pada solstis Juni $\delta = \varepsilon$; pada $\varphi =
\ang{65}$, $\cos h_0 = -\tan\ang{65}\tan\ang{23.446} = -0.9301$, sehingga $h_0 = \ang{158.44} =
\qty{2.7654}{rad}$: Matahari berada di atas ufuk selama $2 \times 158.44/360 \times 24 = 21.1$~jam. Bagian dalam kurung
rumus insolasi adalah $2.7654 \times 0.9063 \times 0.3979 + 0.4226 \times 0.9174 \times 0.3674 =
1.1397$, dan pada jarak rata-rata $\bar Q = 432.90 \times 1.1397 = \qty{493.4}{W.m^{-2}}$. Faktor
jaraknya $(a/r)^2 = 0.9671$, sehingga $\bar Q = \qty{477.1}{W.m^{-2}}$.
\end{example}
% ledger: berger91:rows, berger91:readme, data:b4-orbital

Hukum kedua Kepler menambahkan satu liku. Dekat aphelion Bumi bergerak lebih lambat: musim yang jatuh
di sana lebih lemah tetapi berlangsung lebih lama. Musim panas utara kini sekitar 4.5 hari lebih panjang daripada musim
dingin utara, dan musim semi sekitar tiga hari lebih panjang daripada musim gugur.
% ledger: nasa:milank-detail

\begin{proposition}[Rata-rata tahunan tidak bergantung pada presesi]\label{prop:b2:orbital-forcing:annual}
Energi yang diterima per satuan luas di puncak atmosfer antara dua bujur Matahari yang tertentu
tidak bergantung pada $\varpi$, dan insolasi rata-rata tahunan pada lintang mana pun sama dengan insolasi
orbit lingkaran berjari-jari $a$ dikalikan dengan $1/\sqrt{1 - e^2}$. Di kutub, rata-rata tahunannya adalah
\[
  \bar Q_{\text{kutub}} = \frac{S_0\sin\varepsilon}{\pi\sqrt{1 - e^2}} ,
\]
yang tumbuh seiring kemiringan ekliptika.
\end{proposition}

\begin{proof}
Menurut hukum kedua Kepler, jari-jari menyapu luas dengan laju tetap $\pi ab/T$, dan luas yang tersapu
ketika Bumi berputar sebesar $\mathrm d\lambda$ adalah $\frac12 r^2\,\mathrm d\lambda$, sehingga $\mathrm dt =
T r^2\,\mathrm d\lambda/(2\pi ab)$, dengan $b = a\sqrt{1 - e^2}$. Insolasi pada lintang dan
bujur Matahari tertentu adalah $(a/r)^2$ kali nilainya $Q_a(\lambda)$ pada jarak $a$, sehingga
energi yang diterima adalah $\int Q_a(\lambda)(a/r)^2\,\mathrm dt = \frac{T}{2\pi\sqrt{1 - e^2}}\int
Q_a(\lambda)\,\mathrm d\lambda$: jaraknya saling meniadakan, dan bersamanya $\varpi$. Selama setahun
penuh, $\lambda$ menjalani \ang{360} dan rata-ratanya $\frac{1}{2\pi\sqrt{1-e^2}}\int_0^{2\pi}
Q_a\,\mathrm d\lambda$. Di Kutub Utara $Q_a = S_0\sin\delta = S_0\sin\varepsilon\sin\lambda$
untuk $0 < \lambda < \pi$ dan 0 di luar itu; integralnya $2S_0\sin\varepsilon$.
\end{proof}
% ledger: earth:TSI

Presesi memindahkan sinar Matahari dari satu musim ke musim lain, tidak pernah masuk ke atau keluar dari satu tahun; eksentrisitas
mengubah total tahunan dengan faktor $1/\sqrt{1 - e^2}$, paling banyak \qty{0.16}{\percent} selama sejuta
tahun terakhir; kemiringan sumbu memindahkan sinar Matahari antarlintang, menaikkan rata-rata tahunan di
kutub dan sedikit menurunkannya di khatulistiwa. Hanya pola musiman dan pola lintangnya yang banyak berubah,
dan itulah yang ditunjukkan gambar di bawah.
% ledger: data:b4-orbital, nasa:milank-detail

\begin{omfigure}
\centering
% figdata/bachelor-2/orbital.py part lat
\begin{tikzpicture}
\begin{groupplot}[group style={group size=2 by 1, horizontal sep=1.6cm},
  tick label style={font=\scriptsize}, label style={font=\small}, height=5.4cm,
  width=0.48\linewidth, xmin=-90, xmax=90, xtick={-90,-60,...,90}, xlabel={lintang
  (\unit{\degree})}, ylabel={$\bar Q$ (\unit{W.m^{-2}})}]
\nextgroupplot[ymin=0, ymax=650, title={a\enspace solstis Juni}, title style={font=\scriptsize},
  legend style={font=\scriptsize, at={(0.03,0.97)}, anchor=north west, draw=none},
  legend cell align=left]
\addplot[black, thick] table[x=lat, y=jun_0] {figdata/out/bachelor-2-orbital-lat.dat};
\addlegendentry{masa kini}
\addplot[ocRed, thick, densely dashed] table[x=lat, y=jun_10] {figdata/out/bachelor-2-orbital-lat.dat};
\addlegendentry{\num{10000} tahun yang lalu}
\nextgroupplot[ymin=100, ymax=500, title={b\enspace rata-rata tahunan}, title style={font=\scriptsize},
  legend style={font=\scriptsize, at={(0.5,0.03)}, anchor=south, draw=none},
  legend cell align=left]
\addplot[ocBlue, thick] table[x=lat, y=ann_lo] {figdata/out/bachelor-2-orbital-lat.dat};
\addlegendentry{$\varepsilon = \ang{22.1}$}
\addplot[ocRed, thick, densely dashed] table[x=lat, y=ann_hi] {figdata/out/bachelor-2-orbital-lat.dat};
\addlegendentry{$\varepsilon = \ang{24.5}$}
\end{groupplot}
\end{tikzpicture}
\par\smallskip{\small (a) Insolasi harian rata-rata pada solstis Juni, masa kini dan \num{10000} tahun
yang lalu, ketika musim panas utara jatuh dekat perihelion dan kemiringan sumbu lebih besar: sekitar
\qty{50}{W.m^{-2}} lebih banyak di lintang tinggi utara. (b) Insolasi rata-rata tahunan untuk kedua
kemiringan sumbu ekstrem selama sejuta tahun terakhir, dengan eksentrisitas masa kini: kemiringan yang lebih besar mengirim
lebih banyak sinar Matahari ke kutub dan sedikit lebih sedikit ke daerah tropis. Dihitung dari unsur-unsur orbit
dengan $S_0 = \qty{1360}{W.m^{-2}}$.}
\end{omfigure}
% ledger: data:b4-orbital, data:b4-orbit91

\section{Menghitung kurva insolasi}\label{sec:b2:orbital-forcing:curve}

Kurva insolasi yang dibangun Milankovitch kini merupakan perhitungan rutin. Unsur-unsur orbit
masa lalu dihitung dengan mengintegrasikan gerak planet-planet; ditabulasikan setiap seribu tahun,
unsur-unsur itu menjadi masukannya. Untuk setiap umur, metode di bawah memberikan insolasi pada lintang dan musim mana pun.

\begin{method}[Insolasi masa lalu]\label{met:b2:orbital-forcing:curve}
\begin{enumerate}
  \item Bacalah $e$, $\varpi$, dan $\varepsilon$ untuk umur yang diinginkan.
  \item Pilih musim berdasarkan bujur sejati Matahari $\lambda$: \ang{90} untuk solstis Juni;
        pertengahan Juli dalam tabel-tabel itu sesuai dengan $\lambda = \ang{120}$.
  \item Hitung $\delta$ dari $\sin\delta = \sin\varepsilon\sin\lambda$, dan sudut jam
        terbenamnya Matahari $h_0$ pada lintang $\varphi$ (siang atau malam kutub jika $|\tan\varphi\tan\delta| > 1$).
  \item Hitung faktor jarak $(a/r)^2$ dengan \cref{prop:b2:orbital-forcing:distance}.
  \item Kalikan: $\bar Q = (S_0/\pi)(a/r)^2(h_0\sin\varphi\sin\delta +
        \cos\varphi\cos\delta\sin h_0)$.
\end{enumerate}
\end{method}
% ledger: berger91:readme, data:b4-orbital

Bila diterapkan pada 1001 umur selama sejuta tahun terakhir, metode itu mereproduksi insolasi pertengahan Juli yang ditabulasikan
pada \ang{65} LU sampai dalam \qty{0.012}{W.m^{-2}}, yaitu pembulatan tabel-tabel itu, dan
nilai Juni yang ditabulasikan untuk masa kini secara tepat. Tabel-tabel itu menempatkan titik awal waktu pada 1950.
% ledger: data:b4-orbital, berger91:readme

\begin{omfigure}
\centering
% figdata/bachelor-2/orbital.py part jul65
\begin{tikzpicture}
\begin{groupplot}[group style={group size=1 by 2, vertical sep=0.3cm, xlabels at=edge bottom,
  xticklabels at=edge bottom}, width=0.96\linewidth, xmin=0, xmax=400, xtick={0,50,...,400},
  tick label style={font=\scriptsize}, label style={font=\small}]
\nextgroupplot[height=5.0cm, ymin=370, ymax=590, ylabel={$\bar Q$ pada \ang{65} LU
  (\unit{W.m^{-2}})}, legend style={font=\scriptsize, at={(0.5,1.02)}, anchor=south, draw=none,
  legend columns=3, /tikz/every even column/.append style={column sep=6pt}},
  legend cell align=left]
\addplot[black!45, line width=2.2pt] table[x=age_ka, y=tab] {figdata/out/bachelor-2-orbital-jul65.dat};
\addlegendentry{pertengahan Juli, ditabulasikan}
\addplot[ocRed, thin] table[x=age_ka, y=calc] {figdata/out/bachelor-2-orbital-jul65.dat};
\addlegendentry{pertengahan Juli, dihitung}
\addplot[ocBlue, thin, densely dashed] table[x=age_ka, y=jun65] {figdata/out/bachelor-2-orbital-jul65.dat};
\addlegendentry{solstis Juni, dihitung}
\nextgroupplot[height=2.8cm, ymin=-0.02, ymax=0.02, ytick={-0.01,0,0.01},
  scaled y ticks=false, yticklabel style={/pgf/number format/fixed}, ylabel={residu},
  xlabel={umur (ribuan tahun sebelum 1950)}]
\addplot[only marks, mark=*, mark size=0.35pt, black] table[x=age_ka, y=resid]
  {figdata/out/bachelor-2-orbital-jul65.dat};
\end{groupplot}
\end{tikzpicture}
\par\smallskip{\small Atas: insolasi harian rata-rata pada \ang{65} LU selama \num{400000} tahun terakhir,
pertengahan Juli seperti yang ditabulasikan oleh Berger dan Loutre (abu-abu, lebar) dan seperti yang dihitung dengan metode di atas (merah),
dan pada solstis Juni (biru, putus-putus). Bawah: nilai pertengahan Juli yang dihitung dikurangi nilai yang ditabulasikan, dalam
\unit{W.m^{-2}}.}
\end{omfigure}
% ledger: data:b4-orbit91, data:b4-orbital

Dua ciri menonjol. Kurva itu didominasi oleh \omterm{def:b2:orbital-forcing:elements}{presesi iklim}, dengan puncak-puncak yang berjarak sekitar
\num{20000} tahun yang tingginya membesar dan mengecil mengikuti eksentrisitas; kemiringan sumbu menambahkan
modulasi yang lebih lambat. Dan ayunannya besar: selama \num{800000} tahun terakhir nilai pertengahan Juli
berkisar dari 388 sampai \qty{497}{W.m^{-2}}.
% ledger: data:b4-orbital

\section{Teori Milankovitch dan siklus glasial}\label{sec:b2:orbital-forcing:theory}

\begin{history}[Penyinaran pada enam puluh lima derajat utara]
Milutin Milankovitch, yang lahir pada 1879, membangun teori matematis iklim dari variasi sinar Matahari menurut musim dan lintang. Dengan menggabungkan ketiga variasi orbital,
ia menghitung insolasi di berbagai lintang dan suhu permukaan yang bersesuaian selama
\num{600000} tahun sebelum 1800. Atas saran klimatolog Köppen ia memilih insolasi musim panas
pada \ang{65} LU sebagai besaran yang menentukan: lembaran-lembaran es besar telah tumbuh dekat
lintang itu, dan musim panas yang sejuk mencairkan lebih sedikit salju musim dingin, sehingga salju terakumulasi dari tahun ke
tahun. Ia menganggap kemiringan sumbu sebagai parameter terpenting di antara ketiganya. Selama sekitar lima puluh tahun
teori itu sebagian besar diabaikan; Milankovitch wafat pada 1958. Pada 1976 Hays, Imbrie, dan Shackleton
menemukan dalam sedimen laut dalam dari \num{450000} tahun terakhir puncak-puncak spektral variabilitas iklim
pada sekitar \num{23000}, \num{42000}, dan \num{100000} tahun, yaitu periode presesi, kemiringan sumbu, dan
eksentrisitas, lalu menyimpulkan bahwa perubahan orbit Bumi adalah penyebab mendasar
zaman es Kuarter.
\end{history}
% ledger: eo00:milankovitch, nasa:milank-detail, hays76:periods

Gambar di bawah mengulangi uji mereka pada rekaman-rekaman dari \cref{ch:b2:climate-proxies}. Gambar itu membandingkan
spektrum amplitudo, yaitu kekuatan setiap periode, dari insolasi Juli pada \ang{65} LU dan dari
tumpukan bentik, dalam dua jendela sepanjang \num{700000} tahun.

\begin{omfigure}
\centering
% figdata/bachelor-2/orbital.py part spectrum
\begin{tikzpicture}
\begin{axis}[width=0.9\linewidth, height=5.6cm, xmin=10, xmax=150, ymin=0, ymax=1.08,
  xlabel={periode (ribuan tahun)}, ylabel={amplitudo relatif},
  xtick={19,23,41,100,150}, tick label style={font=\scriptsize}, label style={font=\small},
  legend style={font=\scriptsize, at={(0.5,1.02)}, anchor=south, draw=none, legend columns=3,
  /tikz/every even column/.append style={column sep=6pt}}, legend cell align=left]
\addplot[black!50, thick] table[x=period, y=amp_ins] {figdata/out/bachelor-2-orbital-spectrum.dat};
\addlegendentry{insolasi \ang{65} LU, 0--\qty{700}{ka}}
\addplot[ocRed, thick] table[x=period, y=amp_young] {figdata/out/bachelor-2-orbital-spectrum.dat};
\addlegendentry{LR04, 0--\qty{700}{ka}}
\addplot[ocBlue, thick, densely dashed] table[x=period, y=amp_old] {figdata/out/bachelor-2-orbital-spectrum.dat};
\addlegendentry{LR04, \numrange{1250}{1950}~ka}
\end{axis}
\end{tikzpicture}
\par\smallskip{\small Spektrum amplitudo, masing-masing diskalakan terhadap maksimumnya sendiri, dari insolasi pertengahan Juli
pada \ang{65} LU selama \num{700000} tahun terakhir dan dari tumpukan bentik LR04 pada jendela yang sama
dan pada \num{1250000} sampai \num{1950000} tahun yang lalu (tren dihilangkan, diberi jendela peruncing). Insolasi memiliki
dayanya pada periode presesi, dekat 19 sampai 24 ribu tahun, dan pada periode kemiringan sumbu; volume es
masa kini berpuncak dekat \num{100000} tahun, tempat insolasi hampir tidak memiliki apa-apa; volume es yang lebih tua
berpuncak pada \num{41000} tahun.}
\end{omfigure}
% ledger: data:b4-orbital, data:lr04, hays76:periods

\subsection*{Glasial, interglasial, dan terminasi}

\begin{definition}[Periode glasial; interglasial; terminasi glasial]\label{def:b2:orbital-forcing:glacials}
Dalam rekaman Kuarter, yang disebut \emph{periode glasial}\index{periode glasial} adalah selang waktu ketika lembaran
es besar menutupi benua-benua utara, yang terekam oleh $\delta^{18}$O bentik yang tinggi dan muka laut yang rendah;
yang disebut \emph{interglasial}\index{interglasial} adalah selang waktu dengan sedikit es daratan di Belahan Bumi
Utara di luar Greenland, seperti masa kini. Yang disebut \emph{terminasi
glasial}\index{terminasi glasial} adalah salah satu peralihan besar yang mengakhiri suatu siklus glasial,
dari kondisi glasial penuh ke kondisi interglasial, jauh lebih cepat daripada pertumbuhan yang mendahuluinya.
\end{definition}
% ledger: kohler20:mpt, barker21:terminations, ipcc:ref-periods

Selama setengah juta tahun terakhir periode iklim yang dominan mendekati \num{100000} tahun,
meskipun gaya paksa pada periode itu, yang hanya berasal dari eksentrisitas, dapat diabaikan
dibandingkan dengan gaya paksa kemiringan sumbu dan presesi. Siklus-siklusnya juga timpang: $\delta^{18}$O bentik
merekam penumpukan es bertahap selama sekitar \num{90000} tahun, lalu peluruhan dalam sekitar
\num{10000} tahun, suatu gigi gergaji yang sisi turunnya adalah terminasi. \omterm{def:b2:orbital-forcing:glacials}{Interglasial} terakhir, dari
sekitar \num{129000} sampai \num{116000} tahun yang lalu, berakhir ketika insolasi musim panas pada \ang{65} LU turun
ke salah satu nilai terendahnya dalam \num{800000} tahun terakhir: nilai solstis Juni mencapai
minimumnya, \qty{438.6}{W.m^{-2}}, \num{116000} tahun yang lalu, dan nilai pertengahan Juli minimum terendah keduanya,
\qty{392.2}{W.m^{-2}}, \num{114000} tahun yang lalu. Maksimum Glasial Terakhir datang sekitar
\num{95000} tahun kemudian.
% ledger: barker21:terminations, ipcc:ref-periods, data:b4-orbital, data:b4-orbital-extremes

Seberapa cepat lembaran es dapat mengikuti gaya paksanya? Model paling sederhana memperlakukan volume es $V$ sebagai
besaran yang berelaksasi menuju volume $V_{\mathrm{eq}}$ yang dapat dipertahankan oleh insolasi saat itu, dengan
waktu tanggap $\tau$:
\[
  \frac{\mathrm dV}{\mathrm dt} = \frac{V_{\mathrm{eq}}(t) - V}{\tau} .
\]

\begin{proposition}[Tanggapan yang tertunda]\label{prop:b2:orbital-forcing:lag}
Untuk gaya paksa $V_{\mathrm{eq}} = A\sin\omega t$, penyelesaian model tunda menuju
\[
  V = \frac{A}{\sqrt{1 + \omega^2\tau^2}}\,\sin\bigl(\omega(t - t_{\ell})\bigr),\qquad
  t_{\ell} = \frac{\arctan(\omega\tau)}{\omega} :
\]
es mengikuti gaya paksa lambat hampir sepenuhnya dan gaya paksa cepat hanya dengan lemah, selalu terlambat. Tanggapan
terhadap jumlah sinusoid adalah jumlah tanggapan-tanggapannya: model itu tidak menciptakan periode yang tidak dimiliki
gaya paksanya.
\end{proposition}

\begin{proof}
Cobalah $V = B\sin\omega t + C\cos\omega t$: dengan substitusi, $\omega B\tau = -C$ dan $-\omega C\tau = A -
B$, sehingga $B = A/(1 + \omega^2\tau^2)$, $C = -\omega\tau A/(1 + \omega^2\tau^2)$, amplitudonya
$\sqrt{B^2 + C^2} = A/\sqrt{1 + \omega^2\tau^2}$ dan tunda fasenya $\arctan(\omega\tau)$. Penyelesaian
lain mana pun berbeda dari penyelesaian ini sebesar suatu kelipatan $\mathrm e^{-t/\tau}$, yang lenyap. Persamaan itu
linear, sehingga penyelesaian-penyelesaiannya dapat dijumlahkan.
\end{proof}

\begin{omfigure}
\centering
% figdata/bachelor-2/orbital.py part lag
\begin{tikzpicture}
\begin{axis}[width=0.96\linewidth, height=5.4cm, xmin=0, xmax=400, ymin=-3.4, ymax=4.2,
  xtick={0,50,...,400}, xlabel={umur (ribuan tahun sebelum 1950)},
  ylabel={terbakukan}, tick label style={font=\scriptsize}, label style={font=\small},
  legend style={font=\scriptsize, at={(0.5,1.02)}, anchor=south, draw=none, legend columns=3,
  /tikz/every even column/.append style={column sep=6pt}}, legend cell align=left]
\addplot[black!35, thin] table[x=age_ka, y=forcing] {figdata/out/bachelor-2-orbital-lag.dat};
\addlegendentry{minus insolasi \ang{65} LU}
\addplot[ocRed, thick] table[x=age_ka, y=model] {figdata/out/bachelor-2-orbital-lag.dat};
\addlegendentry{model tunda, $\tau = \qty{10}{ka}$}
\addplot[black, thick] table[x=age_ka, y=lr04] {figdata/out/bachelor-2-orbital-lag.dat};
\addlegendentry{LR04 (es ke atas)}
\end{axis}
\end{tikzpicture}
\par\smallskip{\small Model tunda yang digerakkan oleh insolasi pertengahan Juli pada \ang{65} LU (makin sedikit
penyinaran, makin banyak es), dibandingkan dengan tumpukan bentik, ketiganya dibakukan (rata-rata 0, simpangan baku
1); nilai positif berarti lebih banyak es. Model itu mengikuti liukan presesi tumpukan itu tetapi tidak memiliki
satu pun siklus gigi gergajinya yang besar. Waktu tanggap \qty{10}{ka} adalah pilihan model.}
\end{omfigure}
% ledger: data:b4-orbital, data:lr04

Selama \num{800000} tahun terakhir, model tunda berkorelasi dengan tumpukan itu pada $r = 0.42$, sedangkan
insolasi itu sendiri hanya pada $0.23$: tunda itu membantu, tetapi lembaran es yang linear tidak dapat membuat
siklus \num{100000} tahun dari gaya paksa yang tidak memilikinya. Terminasi pasti melibatkan
perilaku tak linear yang memperkuat dorongan kecil dari insolasi. Kandidatnya adalah umpan balik yang bekerja pada seluruh iklim:
\omterm{def:b2:precambrian-earth:ice-albedo}{umpan balik es--albedo} pada \cref{def:b2:precambrian-earth:ice-albedo}, karena planet dengan lebih banyak
salju dan es memantulkan lebih banyak sinar Matahari; dan karbon dioksida, yang dalam inti es Antarktika naik dalam
skala waktu ribuan tahun setiap kali sirkulasi pembalikan Atlantik melemah, seperti yang
terjadi ketika air lelehan dari lembaran es mengalir ke Atlantik Utara. Di atas suatu ambang karbon
dioksida, kondisi \omterm{def:b2:orbital-forcing:glacials}{interglasial} dapat menjadi stabil.
% ledger: data:b4-orbital, barker21:terminations, eo00:milankovitch

\begin{definition}[Transisi Pleistosen Tengah]\label{def:b2:orbital-forcing:mpt}
Yang disebut \emph{transisi Pleistosen Tengah}\index{transisi Pleistosen Tengah} adalah perubahan, dari sekitar
\num{1250000} sampai \num{700000} tahun yang lalu, dari siklus glasial sekitar \num{41000} tahun dengan
amplitudo kecil ke siklus kuasi-periodik sekitar \num{100000} tahun dengan amplitudo besar, dengan lembaran es
yang lebih tebal, tanpa perubahan yang bersesuaian pada gaya paksa orbital.
\end{definition}
% ledger: clark06:mpt, nasa:milank-detail, kohler20:mpt

Sebelum transisi itu lembaran es mengikuti kemiringan sumbu, seperti yang ditunjukkan spektrum putus-putus pada gambar
di atas; satu hitungan menemukan \omterm{def:b2:orbital-forcing:glacials}{terminasi glasial} pada 67, 88, dan \qty{52}{\percent} siklus
kemiringan sumbu pada Kuarter awal, tengah, dan akhir. Teori orbital saja tidak dapat menjelaskan perubahan itu; kandidatnya adalah dinamika
lembaran es itu sendiri, mungkin setelah tersingkirnya regolit lunak tempat lembaran es itu meluncur,
dan penurunan bertahap karbon dioksida selama minimum glasial. Belum ada penjelasan yang
disepakati secara umum.
% ledger: kohler20:mpt, nasa:milank-detail, clark06:mpt

\section{Karbon dioksida dalam siklus glasial}\label{sec:b2:orbital-forcing:co2}

Inti es Antarktika menunjukkan bahwa karbon dioksida mengikuti siklus glasial: antara 188 dan
\qty{194}{ppm} pada Maksimum Glasial Terakhir, antara 266 dan \qty{282}{ppm} pada \omterm{def:b2:orbital-forcing:glacials}{interglasial} terakhir;
rekaman gabungan pada \cref{ch:b2:climate-proxies} memberikan \qty{189.7}{ppm} untuk 19 sampai
\qty{23}{ka} dan \qty{279.4}{ppm} selama dua ribu tahun sebelum industri. Gaya radiatifnya
tumbuh mendekati logaritma konsentrasinya, \qty{3.93}{W.m^{-2}} untuk setiap penggandaan.
% ledger: ipcc:fig2-34, data:b4-proxies-co2, ipcc:log-co2, ipcc:erf-2xco2

\begin{example}[Orbit lawan karbon dioksida]\label{ex:b2:orbital-forcing:co2}
Dari \qty{189.7}{ppm} ke \qty{279.4}{ppm}, gaya radiatifnya $3.93\log_2(279.4/189.7) = 3.93 \times
0.559 = \qty{2.2}{W.m^{-2}}$, di seluruh bola Bumi dan sepanjang tahun. Orbit, selama sejuta
tahun terakhir, mengubah insolasi rata-rata global dan tahunan paling banyak $340.54 - 340.00 =
\qty{0.54}{W.m^{-2}}$, tetapi insolasi pertengahan Juli pada \ang{65} LU sekitar \qty{109}{W.m^{-2}}.
\end{example}
% ledger: data:b4-proxies-co2, ipcc:erf-2xco2, ipcc:log-co2, data:b4-orbital

Orbit adalah pengatur laju, bukan mesin. Orbit menyebarkan ulang sinar Matahari antara musim dan lintang
dan menentukan waktunya; lembaran es, dengan albedo dan kelambatannya, dan karbon dioksida, dengan
gaya paksa globalnya, mengubah sinyal regional menjadi perubahan iklim global.
% ledger: hays76:periods, barker21:terminations, eo00:milankovitch

\section*{Latihan}\label{sec:b2:orbital-forcing:exercises}

Data untuk latihan, kecuali dinyatakan lain: $S_0 = \qty{1360}{W.m^{-2}}$ (nilai dalam
tabel-tabel itu); tahun tropis \num{365.2422} hari; unsur-unsur orbit dari tabel-tabel itu, yang untuk masa kini
$e = 0.017236$, $\varpi = \ang{101.37}$, $\varepsilon = \ang{23.446}$.
% ledger: berger91:readme, berger91:rows

\begin{exercise}[$\star$]\label{exo:b2:orbital-forcing:1}
Ambil $e = 0.0167$. (a) Hitung perbandingan jarak aphelion $a(1 + e)$ terhadap jarak
perihelion $a(1 - e)$. (b) Nyatakan selisihnya sebagai persentase dari $a$. (c) Dengan faktor berapa
sinar Matahari di perihelion lebih kuat daripada di aphelion?
\end{exercise}
% ledger: earth:ecc, nasa:milank-detail

\begin{exercise}[$\star$]\label{exo:b2:orbital-forcing:2}
Tabel-tabel itu memberikan $e = 0.017236$, $\varpi = \ang{101.37}$ untuk masa kini; $e = 0.019971$, $\varpi =
\ang{293.86}$ untuk \num{10000} tahun yang lalu; $e = 0.019249$, $\varpi = \ang{359.99}$ untuk \num{6000} tahun yang lalu.
(a) Hitung \omterm{def:b2:orbital-forcing:elements}{presesi iklim} untuk setiap umur. (b) Untuk masing-masing, berikan bujur Matahari pada saat
perihelion. (c) Pada umur berapa perihelion paling dekat dengan solstis Juni, dan pada umur berapa dengan
ekuinoks September?
\end{exercise}
% ledger: berger91:rows

\begin{exercise}[$\star$]\label{exo:b2:orbital-forcing:3}
Sumbu rotasi berpresesi dengan periode sekitar \num{25770} tahun relatif terhadap bintang-bintang, sementara
orbit berputar dengan periode sekitar \num{112000} tahun. Anggaplah kedua gerak itu menggeser
perihelion relatif terhadap ekuinoks dalam arah yang sama, sehingga frekuensinya dijumlahkan. (a) Gabungkan
kedua frekuensi itu untuk menemukan periode perihelion
bergerak melalui musim-musim. (b) Bandingkan dengan periode utama \omterm{def:b2:orbital-forcing:elements}{presesi iklim}, \num{19.0},
\num{22.4}, dan \num{23.7} ribu tahun.
\end{exercise}
% ledger: nasa:milank-detail, la04:precession

\begin{exercise}[$\star$]\label{exo:b2:orbital-forcing:4}
Dengan \cref{prop:b2:orbital-forcing:annual} dan $e = 0.0167$, hitung insolasi rata-rata tahunan di
Kutub Utara untuk kemiringan sumbu \ang{22.1} dan \ang{24.5}. Berapa persen insolasi itu bertambah di antara
keduanya?
\end{exercise}
% ledger: nasa:milank-detail, earth:ecc

\begin{exercise}[$\star\star$]\label{exo:b2:orbital-forcing:5}
(a) Hitung secara tepat faktor jarak $(a/r)^2$ pada solstis Juni masa kini dan \num{116000} tahun
yang lalu ($e = 0.043988$, $\varpi = \ang{92.71}$), dan bandingkan dengan pendekatan $1 -
2e\sin\varpi$. (b) Hitung keduanya pada solstis Desember. (c) Mengapa pendekatan itu lebih buruk \num{116000}
tahun yang lalu?
\end{exercise}
% ledger: berger91:rows

\begin{exercise}[$\star\star$]\label{exo:b2:orbital-forcing:6}
Ulangi \cref{ex:b2:orbital-forcing:today} langkah demi langkah: sudut jam terbenamnya Matahari pada \ang{65} LU
pada solstis Juni, panjang siang hari, insolasi pada jarak rata-rata, dan
insolasi masa kini.
\end{exercise}
% ledger: berger91:rows

\begin{exercise}[$\star\star$]\label{exo:b2:orbital-forcing:7}
(a) Dengan memakai $\mathrm dt = T(r/a)^2\,\mathrm d\lambda/(2\pi\sqrt{1 - e^2})$ dan $(r/a)^2 \approx 1 +
2e\cos(\lambda - \varpi)$, tunjukkan bahwa setengah tahun dari ekuinoks Maret sampai ekuinoks September berlangsung
sekitar $T/2 + (2T/\pi)\,e\sin\varpi$. (b) Hitung berapa hari setengah tahun itu melebihi setengah tahun lainnya
pada masa kini. (c) Bandingkan dengan 4.5 hari selisih musim panas utara terhadap musim dingin dan tiga hari selisih
musim semi terhadap musim gugur.
\end{exercise}
% ledger: nasa:milank-detail, berger91:rows

\begin{exercise}[$\star\star$]\label{exo:b2:orbital-forcing:8}
Selama sejuta tahun terakhir eksentrisitas berkisar dari \num{0.0046} sampai \num{0.0565}. (a) Hitung
insolasi rata-rata global dan tahunan $S_0/(4\sqrt{1 - e^2})$ pada kedua ekstrem itu. (b) Seberapa besar, dalam
\unit{W.m^{-2}} dan dalam persen, orbit mengubahnya? (c) Bandingkan dengan rentang
insolasi pertengahan Juli pada \ang{65} LU, 388 sampai \qty{497}{W.m^{-2}}.
\end{exercise}
% ledger: data:b4-orbital

\begin{exercise}[$\star\star$]\label{exo:b2:orbital-forcing:9}
Dengan model tunda dan $\tau = \qty{10}{ka}$, hitung amplitudo relatif dan tunda
volume es untuk gaya paksa sinusoidal dengan periode 23, 41, dan \num{100} ribu tahun. Pita mana yang
diunggulkan lembaran es, dan mengapa?
\end{exercise}

\begin{exercise}[$\star\star\star$]\label{exo:b2:orbital-forcing:10}
Insolasi pertengahan Juli pada \ang{65} LU hampir tidak memiliki daya pada \num{100000} tahun, tetapi volume es
selama \num{700000} tahun terakhir memiliki sebagian besar dayanya di sana. (a) Tunjukkan bahwa jika gaya paksanya $A_1\sin
\omega_1 t + A_2\sin\omega_2 t$, model tunda hanya menanggapi pada $\omega_1$ dan $\omega_2$. (b) Tunjukkan
bahwa tanggapan yang memuat suku $V^2$, misalnya $\mathrm dV/\mathrm dt = (V_{\mathrm{eq}} -
V)/\tau + kV^2$, dapat menciptakan frekuensi $\omega_1 - \omega_2$ dan $\omega_1 + \omega_2$, dengan
menyubstitusikan penyelesaian linear ke dalam $V^2$. (c) Dua suku presesi dengan periode \num{23.7} dan
\num{19.0} ribu tahun berinterferensi dengan periode berapa? Apa kesimpulanmu tentang siklus
\num{100000} tahun?
\end{exercise}
% ledger: la04:precession, barker21:terminations

\begin{exercise}[$\star\star\star$]\label{exo:b2:orbital-forcing:11}
Di Kutub Utara, energi yang diterima per meter persegi selama siang kutub adalah $E = \int S_0(a/r)^2
\sin\delta\,\mathrm dt$ selama setengah tahun ketika $\delta > 0$. (a) Tunjukkan dengan bukti
\cref{prop:b2:orbital-forcing:annual} bahwa $E = S_0T\sin\varepsilon/(\pi\sqrt{1 - e^2})$. (b) Hitung
$E$ dalam \unit{GJ.m^{-2}} untuk $\varepsilon = \ang{22.1}$ dan \ang{24.5}, dengan $e = 0.0167$ dan $T$ =
365.25 hari. (c) Presesi tidak muncul dalam $E$. Mengapa presesi tetap dapat mengubah
pencairan suatu lembaran es?
\end{exercise}
% ledger: nasa:milank-detail

\begin{exercise}[$\star\star\star$]\label{exo:b2:orbital-forcing:12}
(a) Hitung gaya radiatif dari kenaikan karbon dioksida dari \qty{189.7}{ppm} ke
\qty{279.4}{ppm}, dengan \qty{3.93}{W.m^{-2}} per penggandaan dan hukum logaritmik. (b) Rentang
yang diberikan untuk Maksimum Glasial Terakhir dan \omterm{def:b2:orbital-forcing:glacials}{interglasial} terakhir adalah 188 sampai \qty{194}{ppm} dan 266 sampai
\qty{282}{ppm}: hitung gaya radiatif terkecil dan terbesar yang dimungkinkannya. (c) Bandingkan dengan perubahan
insolasi tahunan global pada \cref{exo:b2:orbital-forcing:8} dan dengan ayunan insolasi
musim panas pada \ang{65} LU. Mana yang menjadi pemicu dan mana yang menjadi penguat?
\end{exercise}
% ledger: data:b4-proxies-co2, ipcc:fig2-34, ipcc:erf-2xco2, ipcc:log-co2, data:b4-orbital

\begin{problem}[{Soal akhir pekan --- Sinar Matahari pada \ang{65} LU}]\label{pb:b2:orbital-forcing:1}
\omterm{def:b2:orbital-forcing:glacials}{Interglasial} terakhir berakhir sekitar \num{116000} tahun yang lalu. Berapa banyak penyinaran yang diterima musim panas di
lintang lembaran-lembaran es besar utara saat itu, dibandingkan dengan masa kini? Gunakan data
latihan. Tabel-tabel itu memberikan, untuk \num{116000} tahun yang lalu, $e = 0.0440$, $\varpi = \ang{92.7}$,
$\varepsilon = \ang{22.52}$; untuk masa kini, dibulatkan dengan cara yang sama, $e = 0.0172$, $\varpi = \ang{101.4}$,
$\varepsilon = \ang{23.45}$.

\medskip\noindent\textbf{Bagian I --- Orbit dulu dan kini.}
\begin{enumerate}
  \item Hitung \omterm{def:b2:orbital-forcing:elements}{presesi iklim} dulu dan kini.
  \item Hitung jarak perihelion dan aphelion, dalam satuan $a$, dulu dan kini.
  \item Pada bujur Matahari berapa Bumi berada di perihelion dulu, dan kini? Dekat solstis yang mana?
  \item Dengan faktor berapa sinar Matahari di perihelion lebih kuat daripada di aphelion, dulu dan kini?
  \item Apakah kemiringan sumbu lebih besar atau lebih kecil daripada masa kini? Apa artinya bagi musim panas di
        lintang tinggi?
  \item Hitung insolasi rata-rata global dan tahunan dulu dan kini. Apakah selisihnya berarti?
\end{enumerate}

\medskip\noindent\textbf{Bagian II --- Jarak pada solstis Juni.}
\begin{enumerate}[resume]
  \item Tunjukkan dari \cref{prop:b2:orbital-forcing:distance} bahwa pada solstis Juni $r/a = (1 -
        e^2)/(1 - e\sin\varpi)$.
  \item Hitung $r/a$ pada solstis Juni dulu dan kini.
  \item Hitung $(a/r)^2$ dulu dan kini.
  \item Bandingkan dengan pendekatan $1 - 2e\sin\varpi$.
  \item Berapa persen sinar Matahari solstis lebih lemah daripada pada jarak rata-rata, dulu dan kini?
\end{enumerate}

\medskip\noindent\textbf{Bagian III --- Penyinaran pada \ang{65} LU.}
\begin{enumerate}[resume]
  \item Berikan deklinasi Matahari pada solstis Juni dulu dan kini.
  \item Hitung sudut jam terbenamnya Matahari pada \ang{65} LU pada hari itu, dulu dan kini.
  \item Simpulkan panjang siang harinya.
  \item Hitung insolasi harian rata-rata pada jarak rata-rata.
  \item Hitung insolasi harian rata-rata, dulu dan kini.
  \item Tabel-tabel itu memberikan, untuk pertengahan Juli, 397.60 dan \qty{426.76}{W.m^{-2}}. Mengapa nilai
        solstis lebih tinggi?
  \item Mana yang lebih berpengaruh pada selisih antara dulu dan kini, kemiringan sumbu atau jarak?
\end{enumerate}

\medskip\noindent\textbf{Bagian IV --- Jawaban lembaran es.}
\begin{enumerate}[resume]
  \item Dalam model tunda dengan $\tau = \qty{10}{ka}$, berapa bagian dari perubahan mendadak
        $V_{\mathrm{eq}}$ yang telah dicapai lembaran es setelah 5, 10, dan \qty{20}{ka}?
  \item Insolasi Juni pada \ang{65} LU turun dari maksimum sekitar \num{128000} tahun yang lalu ke
        minimum \num{116000} tahun yang lalu, setengah siklus presesi. Untuk sinusoid dengan periode
        \qty{23}{ka}, berapa bagian gaya paksa yang diikuti lembaran es, dan seberapa terlambat?
  \item Mengapa lembaran es tidak dapat mencapai, dalam satu siklus presesi, ukuran yang dapat dipertahankan oleh insolasi
        musim panas yang rendah?
  \item Es memerlukan sekitar \num{90000} tahun untuk menumpuk sampai Maksimum Glasial Terakhir dan sekitar
        \num{10000} tahun untuk mencair. Mengapa model tunda tidak dapat menghasilkan ketimpangan itu?
  \item Karbon dioksida berada antara 266 dan \qty{282}{ppm} pada \omterm{def:b2:orbital-forcing:glacials}{interglasial} terakhir dan antara
        188 dan \qty{194}{ppm} pada Maksimum Glasial Terakhir. Bagaimana karbon dioksida membantu perubahan
        penyinaran regional menjadi perubahan iklim global?
  \item Sebagai ringkasan: berapa insolasi harian rata-rata pada \ang{65} LU pada solstis Juni ketika
        \omterm{def:b2:orbital-forcing:glacials}{interglasial} terakhir berakhir, \num{116000} tahun yang lalu?
\end{enumerate}
\end{problem}
% ledger: berger91:rows, berger91:readme, data:b4-orbital, ipcc:ref-periods, ipcc:fig2-34, barker21:terminations
